\noindent
{\Large \textbf{\section*{Реферат}}}\par
\addcontentsline{toc}{section}{Реферат} \par

Курсовая работа посвящена предметной области <<Организация выставочной деятельности музея>>. 
Основной проблемой выступает отсутствие единой системы, способной каталогизировать выставки, учитывая перемещение экспонатов внутри музея, а это может привести к потере или порче экспонатов из-за неправильного экспонирования.
В рамках курсовой работы была реализована реляционна база данных, хранящая информацию о сотрудниках музея, выставках, выставочных залах музея, экспонатах и особенностях их хранения и экспонирования.
Для реализации базы данных была использована СУБД PostgreSQL, язык программирования Python и язык управления реляционными БД SQL.
Программу можно использовать для организации выставочной деятельности исторических музеев и музеев изобразительных искусств.

\newpage

\noindent
\vspace{-60pt}

{\Large \textbf{\section*{Введение}}}\par
\addcontentsline{toc}{section}{Введение} \par

Основной проблемой рассматриваемой предметной области «Организация выставочной деятельности музея» выступает отсутствие единой системы, способной осуществлять каталогизацию и контроль перемещения экспонатов внутри музея.
При регулярной организации новых выставок, из-за отсутствия четкого контроля перемещения экспонатов, не исключены потери экспонатов при перевозке из залов в фондохранилища и обратно.

Другая проблема — это сохранность экспонатов. Каждому экспонату нужны свои условия: определенная температура, влажность, свет или специальная витрина. Если поставить старинную картину или вещь в зал, который для этого не предназначен, она может испортиться или вовсе разрушиться. Поэтому важно, чтобы система сама проверяла, подходят ли условия зала для конкретного экспоната.

Для решения указанных проблем была спроектирована и реализована реляционная база данных, хранящая информацию о сотрудниках музея, планируемых и текущих выставках, характеристиках выставочных залов, а также о самих экспонатах и особенностях их хранения и экспонирования.

Для реализации БД была выбрана система управления базами данных PostgreSQL. Данный выбор обусловлен поддержкой возможности масштабирования под большие объёмы данных. Взаимодействие с БД осуществлялось при помощи языка SQL. Логика управления создания и заполнения таблиц была реализована на языке программирования Python, потому что он предоставляет множество библиотек для обработки и визуализации данных, например Mathplotlob.

\newpage